\documentclass[11pt,reqno]{amsart}
\usepackage[T1]{fontenc}
\usepackage{amsmath,amssymb,amsthm,mathtools,booktabs,array,longtable,lmodern}
\usepackage[margin=1in]{geometry}
\usepackage[colorlinks=true,linkcolor=blue,urlcolor=blue,citecolor=blue]{hyperref}
\hypersetup{pdftitle={Cross-polytope packing: corrected candidate and exact reductions}}
\newtheorem{theorem}{Theorem}[section]
\newtheorem{proposition}[theorem]{Proposition}
\newtheorem{lemma}[theorem]{Lemma}
\newtheorem{corollary}[theorem]{Corollary}
\theoremstyle{definition}
\newtheorem{remark}[theorem]{Remark}
\newcommand{\R}{\mathbb R}
\newcommand{\Z}{\mathbb Z}
\newcommand{\Q}{\mathbb Q}
\newcommand{\norm}[1]{\lVert #1\rVert_1}
\newcommand{\vol}{\operatorname{vol}}
\newcommand{\conv}{\operatorname{conv}}
\newcommand{\tr}{\operatorname{tr}}
\newcommand{\sgn}{\operatorname{sgn}}
\newcommand{\rank}{\operatorname{rank}}
\newcommand{\M}{\mathcal M}
\newcommand{\Dold}{D_{\mathrm{old}}}
\newcommand{\Dnew}{D_*}
\newcommand{\C}{C_4}
\title[Cross-polytope packing: corrected candidate and reductions]{Four-dimensional cross-polytope packing:\
a corrected candidate and exact structural reductions}
\author{}
\date{14 September 2026}
\begin{document}
\begin{abstract}
The rational lattice declared optimal in the submitted manuscript is not
locally optimal. We give an explicit rational counterexample, then optimize
its contact configuration exactly. The resulting lattice has parameter
$\zeta^3-6\zeta+3=0$, $1/2<\zeta<2/3$, determinant
$(16-9\zeta+6\zeta^2)/16$, and packing density
$0.8248579173014505\ldots$. We prove that it is a strict local optimum and
the unique optimum in the specified contact region. Global optimality is
not proved.

We also repair the structural argument. A short descent on the coordinates
of a violating vector proves the finite admissibility test for seminorms;
in particular, the alleged singular feasible matrices cannot pass the
correct test. The six classical types and all six test sets are reproduced
by exact integer computations. A complete global reduction is stated that
requires neither the unproved good-basis descent nor parity saturation.
The accompanying Python program verifies the finite arithmetic. It is not
a Lean formalization.
\end{abstract}
\maketitle
\tableofcontents

\section{Status and normalization}
Write
\[
 C_n=\{x\in\R^n:\norm{x}\le1\},\qquad \vol(C_n)=\frac{2^n}{n!}.
\]
A packing lattice for $C_n$ has $\norm{v}\ge2$ for every nonzero lattice
vector. Its determinant is its covolume. In dimension four put
\[
 D_{\mathrm{pack}}=\inf\{\det\Lambda:\Lambda\text{ packs }C_4\},
 \qquad \delta_L(C_4)=\frac{2/3}{D_{\mathrm{pack}}}.
\]
In the conventional terminology of critical determinants,
$D_{\mathrm{pack}}=\Delta(2C_4)=16\Delta(C_4)$.
The manuscript's use of ``the critical determinant of $C_4$'' for
$D_{\mathrm{pack}}$ introduces a factor-of-$16$ error in terminology.
For a fixed body and minimum $2$, arbitrary scaling is not an equivalence
between optimal packing lattices; the relevant ambient symmetries are
signed coordinate permutations.

The rigorous conclusions of this note are
\begin{equation}\label{eq:status}
 \frac23\le D_{\mathrm{pack}}\le \Dnew<\Dold,
 \qquad
 \delta_L(C_4)\ge0.8248579173014505\ldots,
\end{equation}
where
\[
 \Dold=\frac{2514935}{3111696},\qquad
 \Dnew=\frac{16-9\zeta+6\zeta^2}{16},\qquad
 \zeta^3-6\zeta+3=0,\quad \frac12<\zeta<\frac23.
\]
The outstanding global assertion is $D_{\mathrm{pack}}=\Dnew$, followed
by a classification of its equality cases. Neither assertion is supplied
by the submitted manuscript or by this note.

Miller's paper \cite{Miller} supplies the displayed rational matrix and a
numerical search for dense packings; its Section 4.2 does not prove global
optimality. The other Miller reference in the manuscript, arXiv:2501.09245,
concerns kissing numbers \cite{MillerKissing}. It is not a proof of the
packing-density assertion. We make no claim here that the improved
algebraic parameter is absent from all subsequent literature or unpublished
work.

\section{An exact counterexample to the stated main theorem}\label{sec:counter}
Let
\[
 A=\begin{pmatrix}
20&53&-53&-42\\
-62&-22&42&-42\\
-22&-22&-62&62\\
20&53&42&53
\end{pmatrix},\qquad L_{\mathrm{old}}=\Z^4(A/84).
\]
The rows are used throughout as basis vectors. Direct calculation gives
$\det A=-40238960$ and $\det L_{\mathrm{old}}=\Dold$.

\begin{proposition}\label{prop:rational}
The lattice generated by the rows of
\begin{equation}\label{eq:rational}
 B_{\mathrm{rat}}=\frac1{1000}
 \begin{pmatrix}
238&631&-631&-500\\
-738&-262&500&-500\\
-262&-262&-738&738\\
238&631&500&631
\end{pmatrix}
\end{equation}
has minimum $2$ and determinant
\[
 D_{\mathrm{rat}}=\frac{50513750079}{62500000000}
 =0.808220001264.
\]
In particular,
\begin{equation}\label{eq:positive-diff}
 \Dold-D_{\mathrm{rat}}
 =\frac{214636001}{12155062500000000}>0.
\end{equation}
Its packing density is
$125000000000/151541250237=0.8248579169335655\ldots$.
\end{proposition}
\begin{proof}
Let $N=1000B_{\mathrm{rat}}$. Its determinant is $-808220001264$.
For any row vector $v=qN$, duality gives
\begin{equation}\label{eq:inverse-bound}
 |q_j|\le\norm{v}\max_i|(N^{-1})_{ij}|.
\end{equation}
For $\norm{v}\le2000$, the floors of the four right-hand coefficient
bounds are $(1,2,1,1)$. It therefore suffices to check the $134$ nonzero
integer vectors in
$[-1,1]\times[-2,2]\times[-1,1]\times[-1,1]$.
Exact integer evaluation gives minimum $2000$, attained by $30$ vectors.
The same argument for $A$, with radius $168$, gives the same coefficient
box and the same $15$ pairs. The program \texttt{verify.py} performs these
exhaustive evaluations and determinant calculations. Division by $1000$
proves the minimum assertion; the determinant comparison is
\eqref{eq:positive-diff}.
\end{proof}

This counterexample already disproves the main theorem at lines 111--114
of the submitted file. It is independent of every claimed descent or
large computational certificate in that file.

\section{The algebraic candidate and its complete contact analysis}\label{sec:candidate}
Put
\begin{equation}\label{eq:M}
 M(z)=\begin{pmatrix}
2-2z&2+z&-2-z&-2\\
2z-4&-2z&2&-2\\
-2z&-2z&2z-4&4-2z\\
2-2z&2+z&2&2+z
\end{pmatrix},\qquad B(z)=\frac14 M(z).
\end{equation}
Then $B(11/21)=A/84$ and $B(131/250)=B_{\mathrm{rat}}$.

\begin{theorem}\label{thm:family}
For every $z\in[1/2,2/3]$, the lattice $\Z^4B(z)$ has minimum $2$.
Its determinant is
\begin{equation}\label{eq:det-family}
 d(z)=\frac{16-12z+12z^2-z^4}{16}.
\end{equation}
The function $d$ has a unique minimum on this interval, at the unique
root $\zeta$ of $z^3-6z+3$ in the interval. Consequently
\[
 \begin{split}
 \zeta&=0.52397639708186596553084232288\ldots,\\
 \Dnew=d(\zeta)&=0.80822000090353537300824717153\ldots,\\
 \frac{2/3}{\Dnew}&=0.82485791730145054648350034516\ldots.
 \end{split}
\]
\end{theorem}
\begin{proof}
The rows of $M(z)$ have norm $8$ on the stated interval. The exact
admissibility certificate is particularly simple. Generate the type-1
test set of Section~\ref{sec:classification}. For each of its $1156$ pairs
$q$, choose $\tau_q$ to be the coordinate signs of $qA$, assigning $+1$
to a zero coordinate. Exact rational arithmetic verifies
\begin{equation}\label{eq:endpoint-cert}
 \tau_q\cdot qM(1/2)\ge8,\qquad
 \tau_q\cdot qM(2/3)\ge8.
\end{equation}
The scalar function $\tau_q\cdot qM(z)$ is affine. Thus it is at least
$8$ throughout the interval, and $\norm{qM(z)}\ge8$ there. The finite
admissibility theorem, proved in Section~\ref{sec:finite}, establishes
admissibility for all integer vectors, not just the test set.
The program regenerates every $q$ and verifies \eqref{eq:endpoint-cert}.

Direct determinant expansion yields \eqref{eq:det-family}; its numerator
is positive on this interval. Differentiating gives
\[
 d'(z)=-\frac14(z^3-6z+3),\qquad
 d''(z)=\frac34(2-z^2)>0.
\]
The cubic has values $1/8$ and $-19/27$ at the endpoints and is strictly
decreasing between them. This proves existence and uniqueness of $\zeta$
and of the minimum. The relation $\zeta^4=6\zeta^2-3\zeta$ gives
$d(\zeta)=(16-9\zeta+6\zeta^2)/16$.
\end{proof}

Notice that $d'(11/21)=-2/9261<0$. Thus the original rational lattice is
not even stationary in this admissible one-parameter family. The
improvement is small numerically, but exact.

\subsection{All contacts near the optimum}
Table~\ref{tab:contacts} lists the $15$ coefficient pairs, one representative
per pair, their supporting facet signs, and multipliers used below.
For compactness, $0011$ denotes $(0,0,1,1)$, and a sign string denotes a
vector in $\{-1,1\}^4$.

\begin{table}[ht]
\centering
\caption{Contacts and positive multipliers.}\label{tab:contacts}
\begin{tabular}{ccc}
\toprule
$q$&$\sigma_q$&$\lambda_q(z)$\\
\midrule
$0001$&$++++$&$29-56z+26z^2$\\
$0010$&$---+$&$-8+12z+22z^2$\\
$0011$&$-+-+$&$10-8z+2z^2$\\
$0100$&$--+-$&$4+4z+10z^2$\\
$0101$&$-+++$&$24-28z+10z^2$\\
$0110$&$---+$&$14+8z-22z^2$\\
$0111$&$-+++$&$3+4z-4z^2$\\
$1000$&$++--$&$29-56z+26z^2$\\
$1001$&$++-+$&$-9+40z-22z^2$\\
$1010$&$-+-+$&$10-8z+2z^2$\\
$1100$&$-+--$&$24-28z+10z^2$\\
$1101$&$-++-$&$-11+40z-8z^2$\\
$1110$&$-+--$&$3+4z-4z^2$\\
$1111$&$-+-+$&$3$\\
$1211$&$-++-$&$3$\\
\bottomrule
\end{tabular}
\end{table}

Each listed vector has norm $8$ in $\Z^4M(z)$ throughout $[1/2,2/3]$.
There are exactly these $15$ pairs when
\begin{equation}\label{eq:root-interval}
 \frac{5239}{10000}\le z\le\frac{131}{250},
\end{equation}
an interval containing $\zeta$. Here is an exact completeness certificate.
Bound the cubic cofactors and the quartic $-\det M(z)$ on this rational
interval by interval arithmetic. It gives
\[
 8\max_i|(M(z)^{-1})_{ij}|<
 \left(\frac{19}{10},\frac{11}{4},\frac{17}{10},\frac{19}{10}\right)_j.
\]
Hence a vector of norm at most $8$ has coefficient bounds $(1,2,1,1)$.
For the remaining pairs in this box, the same fixed supporting-sign
argument gives norm at least $8+5239/5000$. This last gap is asserted
only for the checked coefficient box; it is not a bound for the entire
second shell. This proves completeness of the contact list. All these
bounds are checked with rational arithmetic in the program.

\subsection{Three variables, followed by a one-variable optimization}
For any real $4\times4$ matrix $Y$, set
\[
 \ell_q(Y)=\sigma_q\cdot(qY).
\]
The $15$ equations $\ell_q(Y)=8$ have rank $13$. Their entire solution
space is
\begin{equation}\label{eq:affine-family}
 Y=M(z)+uU+vV,\qquad z,u,v\in\R,
\end{equation}
where
\[
 U=\begin{pmatrix}0&-1&-1&0\\0&0&0&0\\0&0&0&0\\0&1&0&-1\end{pmatrix},
 \qquad
 V=\begin{pmatrix}1&0&0&1\\0&0&-1&-1\\0&0&0&0\\-1&0&1&0\end{pmatrix}.
\]
This statement follows by row reduction over $\Q$; alternatively one
can verify the displayed solutions and a nonzero $13\times13$ minor
of the contact-equation matrix. The program verifies the rank and the
three independent directions exactly.

\begin{proposition}\label{prop:contact-opt}
Among matrices satisfying all the contact equations and sign inequalities
\[
 \ell_q(Y)=8,\qquad (\sigma_q)_j(qY)_j\ge0
 \quad(q\text{ in Table~\ref{tab:contacts}},\ 1\le j\le4),
\]
the absolute determinant is uniquely minimized by $Y=M(\zeta)$.
\end{proposition}
\begin{proof}
Write $Y$ as in \eqref{eq:affine-family}. The first coordinates of the
$0011$ and $1010$ contacts are $2-4z-v$ and $2-4z+v$, both required
to be nonpositive. Thus $z\ge1/2$. The second coordinates of the
$0111$ and $1110$ contacts are $2-3z+u$ and $2-3z-u$, both required
to be nonnegative. Thus $z\le2/3$.

The decisive identity is
\begin{align}\label{eq:main-identity}
 -\det\bigl(M(z)+uU+vV\bigr)
 ={}&16(16-12z+12z^2-z^4)\notag\\
 &+16(z^2-2z+2)u^2+16z^2uv
       +8(z^2-2z+4)v^2.
\end{align}
It is a polynomial identity, verified by direct expansion. For
$z\in[1/2,2/3]$, the last three terms are at least
$12u^2+20v^2$: the diagonal coefficients are at least $16$ and $24$,
and $|16z^2uv|\le8|uv|\le4u^2+4v^2$.
The first term is positive and has its unique minimum at $\zeta$ by
Theorem~\ref{thm:family}. Hence $\det Y<0$ throughout this region,
and equality at the minimum requires $z=\zeta$, $u=v=0$.
\end{proof}

This closes the optimization of this entire specified contact region.
It does not classify other contact regions.

\subsection{Strict local optimality among all lattices}
\begin{theorem}\label{thm:local}
The lattice $\Z^4B(\zeta)$ is a strict local minimum of the determinant
among packing lattices for $C_4$, in the space of basis matrices.
\end{theorem}
\begin{proof}
It is enough to work with minimum $8$ and $Y_*=M(\zeta)$. Every coordinate
of every contact vector $qY_*$ in the table is nonzero. In a sufficiently
small neighbourhood its sign is therefore unchanged. Any nearby
admissible $Y=Y_*+H$ satisfies
\[
 \ell_q(H)\ge0\quad\text{for every listed }q.
\]
Let $f(Y)=-\det Y$. Cofactor expansion and reduction modulo
$z^3-6z+3$ give the exact identity
\begin{equation}\label{eq:kkt}
 Df(Y_*)[H]=\sum_q\lambda_q(\zeta)\ell_q(H).
\end{equation}
The multipliers are listed in Table~\ref{tab:contacts}. Each is strictly
positive on $[1/2,2/3]$, as is checked by minimizing a quadratic on that
interval. Both the coefficient identity and positivity are verified
exactly by the accompanying program.

On the common kernel of the $\ell_q$, the directions are
$aM'(\zeta)+uU+vV$. The Hessian of $f$ is positive definite there:
\eqref{eq:main-identity} gives a positive quadratic form in $u,v$, and
the $z$ second derivative is $192(2-\zeta^2)>0$. There are no mixed
$z,u$ or $z,v$ terms at $u=v=0$.

For completeness, these facts imply strict local minimality without
invoking an optimization theorem. Use the Frobenius norm for matrix
directions. Otherwise take nonzero feasible
$H_k\to0$ with $f(Y_*+H_k)\le f(Y_*)$. After a subsequence,
$H_k/\|H_k\|_2\to H$ with $\|H\|_2=1$.
The first-order Taylor estimate and \eqref{eq:kkt} imply
$\sum_q\lambda_q(\zeta)\ell_q(H_k)=O(\|H_k\|_2^2)$.
All summands are nonnegative, so $\ell_q(H)=0$ for every $q$.
The second-order Taylor expansion now gives
\[
 \frac{f(Y_*+H_k)-f(Y_*)}{\|H_k\|_2^2}
 =\frac{Df(Y_*)[H_k]}{\|H_k\|_2^2}
   +\frac12D^2f(Y_*)[H,H]+o(1)>0
\]
for large $k$, a contradiction. Dividing matrices by $4$ proves the
assertion for minimum $2$.
\end{proof}

\section{The finite admissibility theorem needs no nonsingularity assumption}\label{sec:finite}
The following argument replaces both the manuscript's singularity
discussion and the unnecessary use of its unproved basis descent in
constructing a complete global search.

Let $p_1,\ldots,p_n\in\Z^n$ be independent rows of an integer matrix
$P$, let $p=|\det P|$, and assume
\[
 S(P)=\conv(\pm p_1,\ldots,\pm p_n),\qquad
 S(P)\cap\Z^n=\{0,\pm p_1,\ldots,\pm p_n\}.
\]
Such a cross-polytope is called clean. Define
$\mu_P(q)=\norm{qP^{-1}}$. Its values on $\Z^n$ are nonnegative
integer multiples of $1/p$.

Define $T(P)$ to consist of all $q\in\Z^n\setminus\{0,\pm p_i\}$
such that, for every $i$ with $(qP^{-1})_i\ne0$, the full-dimensional
cross-polytope obtained by replacing $p_i$ with $q$ is clean.
No condition on a degenerate replacement is needed.

\begin{theorem}[Descent for a violating vector]\label{thm:finite}
Let $F:\R^n\to\R_{\ge0}$ be a seminorm, and suppose $F(p_i)=2$.
If $F(q)\ge2$ for every $q\in T(P)$, then $F(q)\ge2$ for every
nonzero $q\in\Z^n$. Moreover $T(P)$ is finite, and
\begin{equation}\label{eq:universal-bound}
 q\in T(P)\quad\Longrightarrow\quad
 p\,|(qP^{-1})_i|\le n!\quad(1\le i\le n).
\end{equation}
\end{theorem}
\begin{proof}
Suppose a nonzero vector violates admissibility, and choose one, $q$,
with $\mu_P(q)$ minimal. This is possible by the discrete range of
$\mu_P$. Cleanliness implies $\mu_P(q)>1$, since a nonzero lattice
point of $S(P)$ would be a prescribed vertex with $F=2$.

Fix $i$ for which $c_i=(qP^{-1})_i\ne0$. If the replacement were not
clean, it would contain a nonzero integer point $w$ other than its
prescribed vertices, with a representation
\[
 w=\sum_{j\ne i}\alpha_jp_j+\beta q,
 \qquad \sum_{j\ne i}|\alpha_j|+|\beta|\le1.
\]
The possibility $\beta=0$ contradicts cleanliness of $S(P)$. The
possibility $|\beta|=1$ would give $w=\pm q$. Thus $0<|\beta|<1$.
It follows that
\[
 F(w)\le2\sum|\alpha_j|+|\beta|F(q)<2,
\]
while
\[
 \mu_P(w)\le1-|\beta|+|\beta|\mu_P(q)<\mu_P(q).
\]
This contradicts the choice of $q$. All its full-dimensional replacements
are therefore clean, so $q\in T(P)$, another contradiction.

For any clean full-dimensional integer cross-polytope with vertex matrix
$V$, Minkowski's convex body theorem gives
\[
 \frac{2^n}{n!}|\det V|=\vol(S(V))\le2^n,
 \qquad |\det V|\le n!.
\]
For the $i$th replacement, $|\det V|=p|c_i|$. This proves
\eqref{eq:universal-bound}, and hence finiteness.
\end{proof}

\begin{corollary}\label{cor:no-singular}
If a real matrix $B$ satisfies
\[
 \norm{p_iB}=2\quad(1\le i\le n),\qquad
 \norm{qB}\ge2\quad(q\in T(P)),
\]
then $B$ is nonsingular and
$|\det B|\ge2^n/n!$.
\end{corollary}
\begin{proof}
Apply Theorem~\ref{thm:finite} to the seminorm $F(x)=\norm{xB}$.
It proves the infinite admissibility inequalities. If $B$ were singular,
the open symmetric convex set $\{x:\norm{xB}<2\}$ would have infinite
volume: it contains an open cylinder along a nonzero kernel direction.
A bounded symmetric convex subset of that cylinder of volume greater
than $2^n$ would, by Minkowski's theorem for $\Z^n$, contain a nonzero
integer point with $F<2$, a contradiction. For nonsingular $B$, the
volume of this set is $2^n\vol(C_n)/|\det B|$, and Minkowski gives
$|\det B|\ge\vol(C_n)$.
\end{proof}

Thus the manuscript's reported singular matrices passing the complete
finite tests are impossible. The tests regenerated below have exactly
the manuscript's six cardinalities. A claimed singular feasible point
must reflect an incorrect constraint encoding, an incomplete list, or
an inexact feasibility report. The available attachment contains no
such matrix or certificate with which to identify the particular cause.
The determinant dichotomy is true but unnecessarily weak: the singular
alternative can be removed.

\section{Small exact certificates for the classical finite inputs}\label{sec:classification}
\subsection{Clean cross-polytopes as finite subgroups of a torus}
For an integer vertex matrix $P$ with determinant of absolute value $p$,
pass to vertex coordinates:
\[
 \Gamma=\Z^nP^{-1}\supseteq\Z^n,\qquad [\Gamma:\Z^n]=p.
\]
Its quotient $G=\Gamma/\Z^n$ is a finite subgroup of $(\R/\Z)^n$.
Since $p\Gamma\subseteq\Z^n$, represent its elements by
$r\in\{0,\ldots,p-1\}^n$, with coordinates divided by $p$.

\begin{lemma}\label{lem:torus}
The cross-polytope with vertex matrix $P$ is clean if and only if every
nonzero $r\in G$ satisfies
\begin{equation}\label{eq:lee}
 \sum_{j=1}^n\min(r_j,p-r_j)>p.
\end{equation}
\end{lemma}
\begin{proof}
The smallest $\ell^1$ norm of a representative of $r/p+\Z^n$ is
$p^{-1}\sum_j\min(r_j,p-r_j)$. A nonzero coset therefore meets the
closed unit cross-polytope exactly when \eqref{eq:lee} fails. In the
zero coset, its only integer points are $0$ and $\pm e_j$.
\end{proof}

All superlattices $\Gamma\supseteq\Z^4$ of index $p$ can be enumerated
without a cyclicity assumption. Dualize: $\Gamma^*\subseteq\Z^4$ has
index $p$ and a unique upper row Hermite normal form $H$, with positive
diagonal entries $d_1,\ldots,d_4$, product $p$, and
$0\le h_{ij}<d_j$ for $i<j$. Then
$\Gamma=\Z^4(H^{-1})^{\mathsf T}$. Enumerate all such $H$ for
$1\le p\le24$ and apply Lemma~\ref{lem:torus} to the group generated
by the rows of $p(H^{-1})^{\mathsf T}$ modulo $p$.

This requires $201677$ Hermite normal forms. Exactly $54$ survive:
\[
\begin{array}{c|rrrrrr}
p&1&2&3&4&5&6\text{--}24\\\hline
\text{clean superlattices}&1&5&8&16&24&0.
\end{array}
\]
The program compares the complete surviving set, not just these counts,
with the signed permutation orbits of the six cyclic groups generated by
\[
\begin{gathered}
0,\quad \frac{(0,1,1,1)}2,\quad \frac{(1,1,1,1)}2,\quad
\frac{(1,1,1,1)}3,\\
\frac{(1,1,1,2)}4,\quad
\frac{(1,1,2,2)}5.
\end{gathered}
\]
They agree exactly. This is an independently reproducible finite
certificate for the classical six-type classification
\cite{Wolff,Mordell}. It uses only Minkowski's theorem, Hermite normal
form, and the elementary test \eqref{eq:lee}. It does not replace these
mathematical inputs by a numerical search.

\subsection{All six admissibility lists}
Use the vertex normal forms $P_t$ whose first three rows are $e_1,e_2,e_3$
and whose last row is respectively
\[
 (0,0,0,1),\ (0,1,1,2),\ (1,1,1,2),\ (1,1,1,3),\
 (1,1,2,4),\ (1,2,2,5).
\]
The verified classification sharpens \eqref{eq:universal-bound} to
$p|(qP_t^{-1})_i|\le5$. If the last row is $(a_1,a_2,a_3,p)$, write
\[
 r=pqP_t^{-1}=(pq_1-a_1q_4,pq_2-a_2q_4,pq_3-a_3q_4,q_4).
\]
Enumerate $r\in[-5,5]^4$, retain the congruences
$r_i+a_ir_4\equiv0\pmod p$, and test every nonsingular replacement
using \eqref{eq:lee}. Removing the prescribed vertices and identifying
$q$ with $-q$ gives exactly
\begin{equation}\label{eq:testcounts}
 (|T(P_t)|/\{\pm1\})_{t=1}^6=(1156,182,196,65,29,16).
\end{equation}
The actual vectors are included in \texttt{certificate.json} and regenerated
by the checker. Degenerate replacements are not used in this definition
or its proof.

There is an even smaller independent description of type 1. For a
nonnegative coefficient pattern $c$, replacing $e_i$ where $c_i=k>0$
is clean if and only if
\begin{equation}\label{eq:modtest}
 t+\sum_{j\ne i}\min\bigl((tc_j)\bmod k,(-tc_j)\bmod k\bigr)>k
 \quad(1\le t<k).
\end{equation}
Indeed, a nonvertex lattice point in that cross-polytope has coefficient
of $c$ equal to $t/k$, and the other coefficients are minimized by
rounding $tc_j/k$ to integers. Apply \eqref{eq:modtest} to the sorted
nonzero patterns with entries between $0$ and $24$; the bound $24$ uses
only Minkowski, not the six-type classification. Apart from the basis
pattern, precisely the following $14$ survive:
\[
\begin{gathered}
0011,\ 0111,\ 1111,\ 0112,\ 1112,\ 1122,\ 1113,\\
1123,\ 1223,\ 1124,\ 1234,\ 1225,\ 1235,\ 2345.
\end{gathered}
\]
The signed permutation orbits give $1156$ pairs. This is a finite
admissibility certificate; it is not a claim that all minimal vectors
belong to those patterns.

\section{Further critical corrections to the submitted manuscript}\label{sec:audit}
The following are mathematical corrections, in addition to the
counterexample and the false singularity claim already proved.

\subsection{Types are configurations, not mutually exclusive lattice classes}
The assertion at lines 324--330 is false as written, and verification
(V2), at lines 664--673, is false even at the manuscript's proposed
threshold. Let $b_1,\ldots,b_4$ be the rows of $A/84$, and take
\[
 X_1=b_4,\quad X_2=b_3,\quad X_3=b_1,\quad
 X_4=b_1+2b_2+b_3+b_4.
\]
All four vectors are minimal. In the ambient lattice basis
$(b_4,b_3,b_1,b_2)$, they have the prescribed type-3 coordinates
$e_1,e_2,e_3,(1,1,1,2)$. Their cross-polytope is clean: in vertex
coordinates the additional coset is generated by
$(-1,-1,-1,1)/2$, whose minimum representative mass is $2>1$.
Thus $L_{\mathrm{old}}$ itself satisfies the type-3 branch of (V2).

In fact, among the $\binom{15}{4}=1365$ unordered choices of distinct
minimal pairs of $L_{\mathrm{old}}$, the index distribution is
\[
\begin{array}{c|rrrrr}
\text{index (zero means dependent)}&0&1&2&3&4\\\hline
\text{number}&420&831&104&9&1.
\end{array}
\]
Every independent choice is clean; the checker verifies this directly
by the torus criterion. The same contact coefficient set occurs near
$\zeta$. Therefore exclusion of all higher-index configurations is
inappropriate for the improved candidate as well.

The correct implication is only: \emph{if a minimal basis exists, a
type-1 configuration can be chosen}. To make the cases disjoint one
may choose a minimal-vector tuple of least positive determinant, but
then this extra minimality condition must be included in any proposed
exclusion argument. The simpler global reduction below allows overlap.

\subsection{The descent theorem is not established}
The statement at lines 558--563 has no complete proof. More specifically:
\begin{enumerate}
\item The trichotomy does not say that every excess vector belongs to
one of the $14$ test patterns or to the flat family. In its populated
case it produces a smaller minimal vector; this may already have
coordinates in $\{-1,0,1\}^4$, leaving the original excess vector
uncontrolled.
\item A witness swap reducing the coefficients of one vector does not
control the potential summed over \emph{all} minimal vectors. The
proposed LP data involve only selected vectors. The effect on the
remaining minimal vectors requires an additional argument.
\item The plateau distance $\rho$ need not be finite merely because it
is defined as a least path length. Proving that a non-good plateau has
an exit is precisely an omitted part of the descent assertion.
\item The flat family is not covered by the claimed check of the $14$
patterns. A check of finitely many patterns alone does not settle it.
\item The attachment supplies no LP systems, witness lists, Farkas
vectors, or checker from which the claimed $6246480$ certificates can
be reconstructed. We neither use nor regard those counts as a proof.
\end{enumerate}
These are proof gaps, not a counterexample here to the existence of a
good basis. The proposed theorem should be removed from the dependency
chain until it has an independent proof. The global reduction of
Section~\ref{sec:global} does not need it.

\subsection{Parity, swaps, and saturation}
If the new basis vector is $x'_j=\sgn(b_j)y$ with $b_j=\pm1$ and
$b'=\sgn(b_j)b$, then the correct coordinate formula is
\[
 c'_j=c_j,\qquad c'_i=c_i-c_jb'_i\quad(i\ne j).
\]
The formula $c'_j=\sgn(b_j)c_j$ in the draft instead corresponds to
taking $x'_j=y$ without the orientation convention. The absolute-value
potential formula survives, but the coordinate statement needs repair.

The parity floor must be written $\Phi_0(B)$ in general. A change of
basis acts by $\mathrm{GL}_4(\mathbb F_2)$, which does not preserve
Hamming weights. A constant value $32$ follows when every nonzero
parity class has exactly one pair, but not for an arbitrary minimal
configuration.

Two distinct minimal pairs can lie in the same parity class. There is,
however, a useful exact replacement:
\begin{lemma}\label{lem:parity-support}
If $m,n$ are minimal, $m\ne\pm n$, and $m\equiv n\pmod{2\Lambda}$,
then their ambient coordinate supports are disjoint.
\end{lemma}
\begin{proof}
The nonzero lattice vectors $(m+n)/2$ and $(m-n)/2$ have norm at
least $2$ and, by the triangle inequality, at most $2$.
Equality in the first triangle inequality forces $m_jn_j\ge0$ for
every $j$; equality in the second forces $m_jn_j\le0$. Thus
$m_jn_j=0$ for every $j$.
\end{proof}
This argument is also used in \cite{MillerKissing}. In particular,
minimal vectors whose supports all have more than $n/2$ coordinates
occupy distinct nonzero parity classes. It still does not force every
class to occur.

The enumeration $2^6 4^4 8=131072$ at lines 643--654 assumes exactly
one minimal pair in each nonzero parity class. The warning at lines
656--661 correctly recognizes the issue but does not resolve it.
Even after a good basis is proved, missing classes and multiple pairs
in a class must both be handled. There are $40$ pairs of nonzero
ternary vectors in dimension four, not just $15$ possible pairs.

\subsection{Other necessary edits}
\begin{itemize}
\item The open-segment lemma reduces to one line:
$y=tm$, $|t|<1$ implies $\norm{y}=2|t|<2$, so $y=0$.
\item Every lattice basis has a clean vertex cross-polytope, regardless
of minimality. Checking this on $400$ bases adds no evidence.
\item The chain lemma is correct under linear independence. Its
parenthetical dependent extension is false: $x_r=-x_s$, $k=2$ has
minimal endpoints and zero at $t=1$.
\item Correct fractional-generator coordinates are valid and useful:
in the vertex basis the ambient superlattice is generated by $\Z^4$
and $(-a_1,-a_2,-a_3,1)/p$. The problem is an incorrect coordinate
transformation, not fractional generators as such.
\item In an LP, $\norm{qB}\ge2$ is a disjunction of supporting
halfspaces. Fixing signs of four basis vectors does not make all
remaining norm inequalities linear in that same region.
\item The $330$ orbit count for the four vertex sign choices is
consistent: after the first row is made positive, the four columns
form a multiset of four elements from eight sign triples, giving
$\binom{11}{4}=330$. This does not count the additional choices
needed for the admissibility inequalities.
\item Bantegnie's 1965 paper has pages 725--730, not 725--736
\cite{Bantegnie}. Brunngraber's dissertation is \emph{\"Uber Punktgitter}
(Vienna, 1944). The Wolff and Mordell page ranges in the draft are correct.
\end{itemize}

The equivalence between packing and admissibility, primitivity, the
equality case of the norm inequality, the shrinking argument, and index identities
are valid. The trichotomy's individual alternatives and their local
proofs are useful after removing the unsupported descent conclusion.
The precise realization and nonrealization examples reported in the
draft cannot be independently certified from its text alone.

\section{A sound global reduction without good bases or saturation}\label{sec:global}
\subsection{Existence and independent contacts}
Mahler compactness gives a determinant minimizer among packing
lattices. Such a lattice has minimum $2$ and its minimal vectors span
$\R^n$. Here is the missing spanning argument. If their span $W$ is
proper, choose a projection $\pi$ onto a complement and apply
$I-\varepsilon\pi$. The minimal vectors are fixed. The other vectors
have a uniform positive gap above $2$, by discreteness, and
\[
 \norm{(I-\varepsilon\pi)v}
 \ge (1-\varepsilon\|\pi\|_{1\to1})\norm{v}.
\]
For sufficiently small $\varepsilon>0$ this preserves admissibility,
whereas the determinant is multiplied by
$(1-\varepsilon)^{n-\dim W}<1$, a contradiction. See also
\cite[Proposition 4]{Miller}.

Choose an independent tuple of minimal vectors of least positive
determinant. Its cross-polytope is clean: an extra lattice point
inside it is minimal, and replacing a vertex with a nonzero coefficient
strictly between $-1$ and $1$ decreases the determinant. Thus the six
normal forms cover every critical lattice in dimension four.

\subsection{A finite union of compact rational polytopes}
For a type $t$, choose one sign vector $\sigma_i$ for each prescribed
vertex, and one sign vector $\tau_q$ for each $q\in T(P_t)/\{\pm1\}$.
Define a set of real matrices by
\begin{align}\label{eq:polytope}
 \eta\cdot(p_iB)&\le2
       &&(1\le i\le4,\ \eta\in\{-1,1\}^4),\notag\\
 \sigma_i\cdot(p_iB)&=2
       &&(1\le i\le4),\\
 \tau_q\cdot(qB)&\ge2
       &&(q\in T(P_t)/\{\pm1\}).\notag
\end{align}
These are rational linear inequalities and equations. The resulting
polytope is bounded, because the rows of $P_tB$ have norm at most $2$
and $P_t$ is invertible. Every point is nonsingular and admissible by
Corollary~\ref{cor:no-singular}. Conversely, every admissible lattice
in that type's prescribed minimal-vector normal form lies in one of
these polytopes. This is an exact equivalence, including all boundaries.

The determinant has a constant sign on each nonempty polytope, since
the polytope is convex and contains no singular matrix. One rational
feasible point determines the sign. No determinant lower-bound
constraint or special treatment of singular matrices is needed.

\begin{proposition}[Dimension of a minimizing face]\label{prop:face-dim}
Let $\mathcal P$ be a real matrix polytope containing no singular matrix.
Any minimum of $|\det|$ on $\mathcal P$ lies in a face of dimension at
most $n(n-1)/2$.
\end{proposition}
\begin{proof}
Let $B_*$ be a minimizer and $E$ the space of directions of the smallest
face whose relative interior contains $B_*$. Both signs of sufficiently
small displacements in $E$ are feasible. Suppose
$\dim E>n(n-1)/2$. Then the subspace
$B_*^{-1}E=\{B_*^{-1}H:H\in E\}$ meets the space of symmetric matrices
nontrivially, since the latter has dimension $n(n+1)/2$. Write
$H=B_*S\in E$ with $S=S^{\mathsf T}\ne0$.
Along this feasible line,
\[
 |\det(B_*+tH)|=|\det B_*|\det(I+tS)
\]
for small $t$. A two-sided local minimum forces $\tr S=0$, whereas
the second derivative of $\det(I+tS)$ is then
$-\tr(S^2)<0$. This is a contradiction.
\end{proof}

Consequently the remaining scalar polynomial optimizations in
dimension four have degree at most four and at most six variables.
This agrees with the dimensional obstruction already discussed by
Betke--Henk \cite[Section 2, pp. 7--8 of the preprint]{BetkeHenk}.
It is a reduction, not an efficient solution of all resulting problems.
In particular, the determinant need not minimize at a vertex of a
polytope. The cubic parameter in Section~\ref{sec:candidate} exhibits
why solving only rational vertices is inadequate.

\subsection{The precise remaining obligation}
For every nonempty polytope \eqref{eq:polytope}, prove
\begin{equation}\label{eq:remaining}
 s\det B\ge \Dnew,\qquad
 s=\sgn(\det B)\text{ on that polytope}.
\end{equation}
Then classify all equality cases and show that their lattices are
signed-coordinate images of $\Z^4B(\zeta)$, if that uniqueness assertion
is true. One must allow the same lattice to occur in several types.

For feasibility pruning, rational Farkas certificates apply directly
to the linear systems. For determinant bounds, a separate polynomial
certificate is required; Farkas certificates alone do not establish
\eqref{eq:remaining}. Candidate critical faces can be further rejected
by the first derivative or by a negative second derivative along a
feasible tangent direction. The explicit rank, multiplier, and
positive-quadratic-form calculation above is a model of the kind of
short final certificate worth seeking for each surviving contact region.

This finite cover is complete and avoids both disputed hypotheses in
the draft. Its exhaustive global evaluation has not been performed.
In particular, this note does not describe \eqref{eq:remaining} as a
routine last check or as an already certified theorem.

\section{What extends to higher dimensions}\label{sec:higher}
The following results hold in every dimension, with the proofs already
given or with the indicated short argument.
\begin{enumerate}
\item The violating-vector descent, including its application to
seminorms, works for any norm precomposed with a linear map. For a
clean integer frame $P$, all necessary test vectors satisfy
$p|(qP^{-1})_i|\le n!$. It is not a specifically four-dimensional
phenomenon.
\item Every clean cross-polytope has index at most $n!$. The torus
criterion \eqref{eq:lee} and enumeration by dual Hermite normal forms
are valid in every dimension. They allow noncyclic quotients from the
start. There is no basis for importing the six-type list into dimension
five or higher.
\item For the cross-polytope, two distinct minimal pairs in one
parity class have disjoint supports. Thus a collection all of whose
vectors have support size greater than $n/2$ has at most $2^n-1$
pairs. Neither saturation nor a good basis follows.
\item A critical lattice has spanning minimal vectors and a clean
minimal frame. The compact polytope construction applies to every
rational centrally symmetric polytope. The dimension bound for a
minimizing face is $n(n-1)/2$: it is $10$ in dimension five and $15$
in dimension six, with determinant degrees $5$ and $6$ respectively.
\item A useful equivalent local test is this. Select one supporting
functional for each contact. The rank of the resulting linear
contact equations at a local minimum is at least $n(n+1)/2$.
Indeed, their common kernel admits two-sided feasible displacements
locally: the chosen functionals keep the contacts outside the body,
and the inactive vectors have a uniform gap. Apply the symmetric
matrix argument of Proposition~\ref{prop:face-dim}. In particular
there are at least $n(n+1)/2$ contact pairs. The contact-number bound
is classical \cite{SwinnertonDyer}.
\item If one instead chooses an independent tuple of minimal vectors
of \emph{largest} determinant, every minimal vector has coordinates
of absolute value at most $1$ in that frame, by the replacement
determinant formula. The frame may have index greater than $1$;
its coordinates lie in $p^{-1}\Z^n$, not necessarily in $\Z^n$.
Also $p\le n!$: the frame determinant is at most $2^n$ by Hadamard,
whereas the packing lattice determinant is at least $2^n/n!$.
This gives a sound rational-coordinate alternative to the unproved
good-basis theorem, but may substantially enlarge the list of indices.
\item For a rational centrally symmetric polytope, these finite
descriptions also imply that a critical determinant is a real
algebraic number and that an algebraic critical basis can be chosen.
The minimum is attained; its graph and minimizing set are
semialgebraic over $\Q$, so real closed field transfer gives a
minimizer over the real algebraic numbers. This is an existence
statement, not a practical complexity bound. On a nonconstant
one-dimensional contact family, an interior critical parameter
satisfies a polynomial of degree at most $n-1$.
\end{enumerate}

The factorial index bound and increasing face dimension explain the
computational difficulty in higher dimensions. The realistic extension
is to close selected contact families by polynomial identities and
positive multipliers, while retaining the full finite cover as a
correctness specification. No global density formula in dimensions
five or higher follows from the present work.

\section{Exact verification and the appropriate Lean targets}\label{sec:lean}
The supplied archive contains \texttt{verify.py}, this task's
\texttt{certificate.json}, and a short README. The checker uses only
Python's standard library. All accept/reject decisions use integer or
rational arithmetic. Decimal arithmetic is used only to print
approximations after the exact checks have passed.

\begin{center}
\begin{tabular}{p{0.56\linewidth}p{0.32\linewidth}}
\toprule
Finite item & Exact outcome\\\midrule
Original and rational counterexample: complete short-vector check
& $134$ nonzero coefficient vectors each; $15$ minimal pairs\\
Clean cross-polytope classification
& $201677$ HNFs; $54$ survivors; six orbits\\
Six admissibility sets
& $1156,182,196,65,29,16$ pairs\\
Independent type-1 modular check
& Exactly the $14$ listed patterns\\
Admissibility of the entire interval family
& $1156$ fixed supporting functionals at two rational endpoints\\
Contact affine space
& Rank $13$, three explicit directions\\
Determinant and multiplier certificates
& Polynomial identities and positive quadratics\\
Global inequality \eqref{eq:remaining}
& Not established\\
Lean kernel verification
& Not performed\\\bottomrule
\end{tabular}
\end{center}

The first Lean milestone should be the rational counterexample, not
the false proposed optimum. It needs an inverse/cofactor coefficient
bound and a small finite norm check. Next formalize the polynomial
identity \eqref{eq:main-identity}, the cubic root specified by an
isolating interval, the positive-multiplier identity, and the local
minimum argument. These statements have explicit correct targets.

For the global project, formalize Theorem~\ref{thm:finite} and
Corollary~\ref{cor:no-singular} abstractly, then the torus criterion and
the complete finite cover. One may initially import the classical
six-type theorem as a clearly identified mathematical dependency; a
fully self-contained development must also verify the HNF enumeration
or supply a formal proof of that classification. The same applies to
Minkowski's theorem and compactness if they are not already available
in the chosen library version.

The Python checker is reproducible evidence, not a Lean proof. The
global theorem should remain an explicitly unproved target until
certificates for \eqref{eq:remaining} and its equality cases actually
exist. There is no reason to formalize the draft's unsupported descent
or its impossible singular-feasibility reports.

\begin{thebibliography}{9}
\bibitem{Miller}
N. Miller, \emph{A design criterion for the Rayleigh fading wiretap
channel based on $\ell^1$-norm theta functions},
\href{https://arxiv.org/abs/2405.04143}{arXiv:2405.04143} (2024),
especially Section 4 and Table 1. Published version:
\href{https://doi.org/10.1137/24M1658784}{doi:10.1137/24M1658784} (2025).
\bibitem{MillerKissing}
N. Miller, \emph{On the kissing number of the cross-polytope},
\href{https://arxiv.org/abs/2501.09245}{arXiv:2501.09245} (2025), Section 3.
\bibitem{Wolff}
K. H. Wolff, \emph{\"Uber kritische Gitter im vierdimensionalen Raum
$(R_4)$}, Monatsh. Math. \textbf{58} (1954), 38--56,
\href{https://doi.org/10.1007/BF01478562}{doi:10.1007/BF01478562}.
The present note supplies its own finite-test proof; the bibliographic
record, not the full original article, was accessible during this audit.
\bibitem{Mordell}
L. J. Mordell, \emph{Lattice octahedra}, Canad. J. Math. \textbf{12}
(1960), 297--302,
\href{https://doi.org/10.4153/CJM-1960-025-4}{doi:10.4153/CJM-1960-025-4}.
\bibitem{Bantegnie}
R. Bantegnie, \emph{Sur l'indice de certains r\'eseaux de $\R^4$ permis
pour un octa\`edre}, Canad. J. Math. \textbf{17} (1965), 725--730.
Bibliographic details are also recorded in \cite{BetkeHenk}.
\bibitem{BetkeHenk}
U. Betke and M. Henk, \emph{Densest lattice packings of 3-polytopes},
Comput. Geom. \textbf{16} (2000), 157--186;
\href{https://arxiv.org/abs/math/9909172}{arXiv:math/9909172}.
\bibitem{SwinnertonDyer}
H. P. F. Swinnerton-Dyer, \emph{Extremal lattices of convex bodies},
Proc. Cambridge Philos. Soc. \textbf{49} (1953), 161--162,
\href{https://doi.org/10.1017/S0305004100028188}{doi:10.1017/S0305004100028188}.
\bibitem{Cassels}
J. W. S. Cassels, \emph{An introduction to the geometry of numbers},
Springer, 1959.
\bibitem{GruberLekkerkerker}
P. M. Gruber and C. G. Lekkerkerker, \emph{Geometry of numbers},
second edition, North-Holland, 1987.
\end{thebibliography}
\end{document}
